% =====================================================================
%  CHAPITRES 1 ET 2 — Comment on devient bilingue + le plan de 12 mois
% =====================================================================
\chapter{Comment on devient vraiment capable de parler}

\begin{quote}
\itshape
On ne devient pas bilingue en \textbf{comprenant} une langue. On devient bilingue en
\textbf{produisant} une langue, beaucoup, souvent, sans attendre d'être parfait.
\end{quote}

Avant de toucher à la grammaire, il faut comprendre comment le cerveau fabrique une langue.
Cinq minutes de lecture ici vous feront gagner des mois.

\section{Ce que dit la recherche (en quatre idées simples)}

\subsection{Idée 1 --- L'input compréhensible : comprendre à 90\,\%}
Le linguiste Stephen Krashen a montré qu'on acquiert une langue quand on comprend des
messages \emph{légèrement} au-dessus de son niveau. Ni trop facile (on n'apprend rien),
ni trop difficile (on décroche). La règle pratique :

\begin{memo}[La règle des 90\,\%]
Choisissez un texte ou un audio où vous comprenez \textbf{9 mots sur 10}. Si vous comprenez
moins de 8 mots sur 10, c'est trop dur : descendez d'un niveau. Si vous comprenez
10 mots sur 10, c'est trop facile : montez.
\end{memo}

\subsection{Idée 2 --- L'output : parler n'est pas la conséquence de l'apprentissage, c'en est le moteur}
Pendant longtemps on a cru : « d'abord je stocke, ensuite je parle ». Faux. Parler
\emph{fabrique} la compétence : quand vous cherchez un mot et ne le trouvez pas, votre
cerveau enregistre un « trou » qu'il comblera durablement la prochaine fois qu'il rencontrera
le mot. C'est ce qu'on appelle l'\emph{effet de production} et l'\emph{hypothèse de
l'output} (Merrill Swain). Conséquence pour vous : \textbf{parlez dès la première semaine},
même mal, même seul.

\subsection{Idée 3 --- La répétition espacée : revoir juste avant d'oublier}
Vous oubliez 70\,\% d'un mot nouveau en 24 heures si vous ne le revoyez pas. Mais si vous le
revoyez à J+1, J+3, J+7, J+15, J+30, il reste. C'est le principe des logiciels de
répétition espacée (\en{spaced repetition systems}, SRS) comme \textbf{Anki}. Un mot appris
sans SRS coûte 10 fois plus cher en temps qu'un mot appris avec SRS.

\subsection{Idée 4 --- Le sommeil et la régularité battent l'intensité}
Deux heures par jour pendant un an (730 h) valent infiniment mieux que dix heures par jour
pendant trois semaines. La consolidation de la mémoire se fait la nuit ; le cerveau a besoin
de contacts \emph{répétés}. Visez la \textbf{chaîne ininterrompue} : 365 jours d'affilée
valent mieux que 250 jours parfaits.

\section{Combien d'heures faut-il ? (le calcul honnête)}
Le Cadre européen commun de référence pour les langues (CECRL) et les estimations de
l'institut américain FSI donnent des ordres de grandeur fiables pour un francophone :

\begin{tableau}[Heures d'étude pour un francophone --- anglais]
\begin{tabularx}{\linewidth}{@{}p{2.2cm}p{3.4cm}X@{}}
\entete{Niveau} & \entete{Volume cumulé} & \entete{Ce que vous savez faire} \tabularnewline
\hline
\textbf{A1} & 80 à 120 h & Se présenter, formules de survie, phrases très simples. \tabularnewline
\textbf{A2} & 200 à 300 h & Échanges courants : achats, famille, travail, trajets, passé et futur simples. \tabularnewline
\textbf{B1} & 400 à 550 h & Se débrouiller en voyage, raconter, expliquer, donner un avis, suivre une série avec sous-titres anglais. \tabularnewline
\textbf{B2} & 700 à 900 h & Parler avec aisance et spontanéité, argumenter, comprendre les informations et la plupart des films. \textbf{C'est l'objectif de ce livre.} \tabularnewline
\textbf{C1} & 1\,000 à 1\,400 h & Nuancer, plaisanter, rédiger des textes complexes, travailler en anglais sans effort. \tabularnewline
\end{tabularx}
\end{tableau}

\begin{astuce}[Votre budget temps]
90 minutes par jour $\times$ 6 jours $\times$ 52 semaines = \textbf{468 heures d'étude active}.
Ajoutez l'immersion passive (musique, séries, téléphone en anglais, radio) : environ
\textbf{2 heures par jour}, soit 730 heures de contact. Total annuel : plus de
\textbf{1\,100 heures d'exposition}. Avec un départ A1/A2, cela vous place \textbf{entre B2
et C1}. Le chiffre n'est pas une promesse : c'est la seule chose que vous contrôlez.
\end{astuce}

\section{Les cinq systèmes à faire tourner en parallèle}
Une langue n'est pas une matière scolaire, c'est cinq compétences qui progressent à des
vitesses différentes. Voici les cinq « moteurs » et la part de votre temps quotidien à leur
consacrer.

\begin{tableau}[Les 5 moteurs et leur poids dans votre journée de 90 min]
\begin{tabularx}{\linewidth}{@{}p{3.6cm}p{2.0cm}X@{}}
\entete{Moteur} & \entete{Temps} & \entete{Activités concrètes} \tabularnewline
\hline
\textbf{1. Oreille} (compréhension orale) & 25 min &
Écoute intensive (5 passes sur 2 min), dictée, podcasts, séries, chansons. \tabularnewline
\textbf{2. Bouche} (production orale) & 25 min &
Shadowing, lecture à voix haute, auto-conversation, dialogues du livre, appels avec
partenaire. \tabularnewline
\textbf{3. Vocabulaire} & 15 min &
Anki (révisions + 15 mots nouveaux), blocs de phrases, familles de mots. \tabularnewline
\textbf{4. Grammaire / lecture} & 15 min &
Une unité de ce livre + 5 exercices + un texte court. \tabularnewline
\textbf{5. Écriture} & 10 min &
Journal : 5 à 10 phrases sur votre journée, puis correction. \tabularnewline
\end{tabularx}
\end{tableau}

\begin{piege}[Les francophones négligent le moteur n°2]
L'école africaine et française enseigne l'anglais \emph{par l'écrit} : on lit, on conjugue,
on fait des thèmes. Résultat : des adultes qui savent accorder un \en{present perfect} mais
qui restent muets. Ce livre inverse la proportion : \textbf{un tiers de votre temps doit
sortir de votre bouche}, chaque jour, dès le premier jour.
\end{piege}

\section{Les six mythes qui font perdre du temps}
\begin{enumerate}
  \item \textbf{« Je suis nul en langues.} » Faux : vous avez appris le français, le wolof,
  ou une autre langue sans manuel, en 3 ans, en écoutant et en osant. Vous savez faire.
  \item \textbf{« Il faut un professeur particulier.} » Utile, pas indispensable. Un
  partenaire de conversation gratuit (application, club, église, camarade) 2 fois par
  semaine suffit si vous travaillez sérieusement seul le reste du temps.
  \item \textbf{« Il faut vivre dans un pays anglophone.} » L'immersion aide énormément,
  mais on peut fabriquer une immersion artificielle chez soi (chapitre 4). Des millions de
  personnes ont atteint C1 sans quitter leur pays.
  \item \textbf{« Je dois connaître tous les mots.} » 3\,000 mots couvrent 95\,\% d'une
  conversation quotidienne. 1\,000 mots en couvrent 85\,\%. La quantité ne sert à rien sans
  les mots \emph{utiles}.
  \item \textbf{« Je parlerai quand je serai prêt.} » On n'est jamais prêt. On parle, donc
  on devient prêt.
  \item \textbf{« L'accent parfait ou rien.} » Votre accent français (ou sénégalais) n'est
  pas un défaut : c'est votre identité. L'objectif est l'\textbf{intelligibilité} --- qu'on
  vous comprenne du premier coup --- pas l'imitation d'un acteur de Hollywood.
\end{enumerate}

\section{Les trois leviers de la motivation sur douze mois}
\begin{itemize}
  \item \textbf{Le levier visuel} : imprimez l'annexe F (journal de progression) et collez-la
  au mur. Une case cochée par jour. Une chaîne de 365 cases est psychologiquement
  imparable : on ne veut pas la casser.
  \item \textbf{Le levier social} : annoncez votre objectif à deux personnes. Mieux :
  trouvez un partenaire d'apprentissage (même à distance) et fixez un rendez-vous hebdomadaire.
  \item \textbf{Le levier de la preuve} : enregistrez-vous le jour 1 (\en{``My name is\ldots''})
  et conservez le fichier. Réenregistrez-vous chaque mois. Écouter sa propre voix à trois
  mois d'intervalle est la meilleure motivation qui existe.
\end{itemize}

\begin{mission}[Jour 0 --- aujourd'hui, 20 minutes]
\begin{enumerate}
  \item Enregistrez-vous sur votre téléphone en anglais pendant 2 minutes : présentez-vous,
  dites ce que vous faites dans la vie, ce que vous voulez dans un an. Nommez le fichier
  \en{jour-000.m4a}. Ne l'effacez jamais.
  \item Écrivez en dessous de l'enregistrement, en français, \textbf{pourquoi} vous voulez
  parler anglais (travail, religion, études, famille, voyages, Internet). Soyez précis.
  Relisez cette phrase chaque lundi.
  \item Cochez la première case de votre journal de progression.
\end{enumerate}
\end{mission}


% =====================================================================
\chapter{Le plan des 12 mois, mois par mois}

Ce chapitre est le c\oe ur opérationnel du livre. Il vous dit \textbf{quoi faire, quand,
combien de temps et comment vérifier}. Tout le reste du livre est au service de ce plan.

\section{La journée type (90 minutes)}
Voici le squelette quotidien. Adaptez les horaires, gardez les proportions.

\begin{tableau}[Votre séance quotidienne --- 90 minutes]
\begin{tabularx}{\linewidth}{@{}p{1.6cm}p{4.4cm}X@{}}
\entete{Ordre} & \entete{Bloc} & \entete{Détail} \tabularnewline
\hline
1 & \textbf{Échauffement} (5 min) &
Lecture à voix haute d'un paragraphe déjà connu (du chapitre en cours ou d'un dialogue
appris). Objectif : remettre la bouche en mouvement, sans réfléchir. \tabularnewline
2 & \textbf{Écoute intensive} (20 min) &
Un audio court (1 à 3 min) travaillé en 5 passes : écoute globale \ra écoute avec le texte
\ra écoute sans le texte \ra dictée d'une phrase \ra shadowing. Protocole complet au
chapitre 16. \tabularnewline
3 & \textbf{Vocabulaire SRS} (15 min) &
Révisions Anki du jour + 15 mots/blocs nouveaux issus de l'unité du livre. Toujours à voix
haute, toujours en phrase. \tabularnewline
4 & \textbf{Unité du livre} (25 min) &
Lecture de l'unité (règles, tableaux, pièges) + exercices écrits + vérification des corrigés. \tabularnewline
5 & \textbf{Production orale} (15 min) &
La « mission orale » de l'unité : répéter 10 phrases, se raconter sa journée, rejouer le
dialogue en changeant les détails, s'enregistrer. \tabularnewline
6 & \textbf{Journal écrit} (10 min) &
5 à 10 phrases sur votre journée, en anglais, sans dictionnaire d'abord, puis correction
avec un dictionnaire ou un correcteur. \tabularnewline
\end{tabularx}
\end{tableau}

\begin{astuce}[Vous n'avez que 30 minutes ?]
Ne supprimez pas au hasard. Gardez toujours, dans l'ordre : \textbf{(1)} les 15 min de
vocabulaire SRS, \textbf{(2)} 10 min de parole à voix haute, \textbf{(3)} 5 min d'écoute.
La grammaire peut attendre un jour ; la répétition espacée et la parole, jamais. Une
journée à 30 minutes vaut mieux qu'une journée à zéro.
\end{astuce}

\begin{astuce}[Le temps mort devient du temps d'anglais]
Vous avez chaque jour 1 à 2 heures de temps mort (transport, cuisine, marche, attente,
ménage). Transformez-les en immersion : écouteurs, un podcast ou une playlist. Cela double
votre exposition sans effort supplémentaire. C'est la différence entre ceux qui réussissent
en 12 mois et ceux qui échouent en 3 ans.
\end{astuce}

\section{La semaine type}
\begin{tableau}[Semaine de 6 jours + 1 jour de bilan]
\begin{tabularx}{\linewidth}{@{}p{2.0cm}p{4.0cm}X@{}}
\entete{Jour} & \entete{Accent principal} & \entete{Contenu} \tabularnewline
\hline
\textbf{Lundi} & Grammaire neuve & Séance complète ; nouvelle unité du livre. Fixez l'objectif de la semaine. \tabularnewline
\textbf{Mardi} & Vocabulaire + écoute & Séance complète ; 1\textsuperscript{re} révision de l'unité de lundi à voix haute (5 min). \tabularnewline
\textbf{Mercredi} & \textbf{Parole} & Séance + 30 min de conversation (partenaire, application, ou monologue enregistré sur un sujet du livre). \tabularnewline
\textbf{Jeudi} & Grammaire + lecture & Séance complète ; lecture extensive 15 min (texte facile, plaisir). \tabularnewline
\textbf{Vendredi} & Vocabulaire + écoute & Séance + dictée de 5 phrases + chanson travaillée. \tabularnewline
\textbf{Samedi} & \textbf{Intégration} & Séance allégée + 45 min de « projet » : série, appel, présentation, écriture d'un texte long. \tabularnewline
\textbf{Dimanche} & Bilan (45 min) & Test rapide de la semaine, relecture des corrigés, préparation du matériel de la semaine, enregistrement vocal de 3 min. \tabularnewline
\end{tabularx}
\end{tableau}

\section{Les douze mois en détail}

\subsection{Mois 1 --- Démarrer fort (A1)}
\textbf{Priorités :} les sons (chapitre 5), les 300 mots les plus fréquents, \en{be},
\en{have}, le présent simple, se présenter, les nombres, l'alphabet.

\begin{itemize}
  \item \textbf{Chapitres du livre :} 1 et 2 (la méthode, le plan), 3 et 4 (les huit techniques, les outils), 5 (les 44 sons), 8 (unités 1 à 4 : \en{be}, \en{have}, le présent simple, les articles), 12 (thèmes 1 à 4).
  \item \textbf{Vocabulaire :} 15 mots/jour $\times$ 26 jours = \textbf{390 mots}.
  \item \textbf{Écoute :} 15 min/jour d'un podcast pour débutants (voir annexe E), vitesse 0,75×.
  \item \textbf{Parole :} chaque jour, 10 phrases à voix haute + présentation de soi
  (2 min) une fois par semaine, enregistrée.
\end{itemize}
\textbf{Jalon de fin de mois :} vous vous présentez pendant 2 minutes sans lire ;
vous comprenez et dites 200 phrases de survie ; vous reconnaissez les 44 sons de l'anglais.

\subsection{Mois 2 --- Construire la phrase (A1)}
\textbf{Priorités :} poser des questions, nier, les articles, le pluriel, \en{there is/are},
les prépositions de lieu et de temps, les possessifs, décrire.

\begin{itemize}
  \item \textbf{Chapitres :} 6 (accent, rythme, \en{connected speech}), 8 (unités 5 à 10 : questions, négation, pluriels, \en{there is/are}, prépositions, possessifs), 12 (thèmes 5 à 8), 14 (situations 1 à 8).
  \item \textbf{Vocabulaire :} 600 mots cumulés.
  \item \textbf{Parole :} décrire une pièce, une personne, sa famille, son quartier (3 min chacun).
\end{itemize}
\textbf{Jalon :} vous posez 30 questions types sans hésiter ; vous décrivez une image
pendant 2 minutes.

\subsection{Mois 3 --- Le passé et le futur (A1\ra A2)}
\textbf{Priorités :} \en{past simple} (et les 60 verbes irréguliers essentiels),
\en{past continuous}, \en{going to}, \en{will}, comparatifs, raconter.

\begin{itemize}
  \item \textbf{Chapitres :} 9 (unités 1 à 4 : \en{past simple}, \en{past continuous},
  \en{going to}, \en{will}), annexe A (verbes irréguliers, groupes 1 et 2),
  12 (thèmes 9 à 12), 14 (situations 9 à 14).
  \item \textbf{Vocabulaire :} 900 mots cumulés.
  \item \textbf{Parole :} raconter sa journée d'hier, un souvenir d'enfance, ses projets
  de l'année (3 min chacun).
\end{itemize}
\textbf{Jalon :} vous racontez une histoire de 2 minutes au passé ; vous connaissez
60 verbes irréguliers par c\oe ur.

\subsection{Mois 4 --- Le mur du \en{present perfect} (A2)}
\textbf{Priorités :} \en{present perfect} contre \en{past simple} (le plus gros piège
francophone), \en{for}/\en{since}, \en{just}/\en{already}/\en{yet}, \en{used to},
modaux de base.

\begin{itemize}
  \item \textbf{Chapitres :} 9 (unités 5 à 8 : \en{present perfect}, \en{for}/\en{since}, \en{used to}, modaux de base), 18 (erreurs 1 à 30), 14 (situations 15 à 20).
  \item \textbf{Vocabulaire :} 1\,200 mots cumulés --- santé, ville, transports, courses.
  \item \textbf{Parole :} parler de son expérience (\en{I have never\ldots}, \en{I have been\ldots}),
  situations de voyage.
\end{itemize}
\textbf{Jalon :} dans un test de 20 phrases, vous choisissez correctement \en{past simple}
ou \en{present perfect} dans 18 cas sur 20.

\subsection{Mois 5 --- Consolider et travailler (A2)}
\textbf{Priorités :} révision de tous les temps, voix passive simple, \en{should}/\en{must}/
\en{have to}, vocabulaire du travail et des métiers, premier roman simplifié.

\begin{itemize}
  \item \textbf{Chapitres :} 10 (unités 1 à 3 : passif simple, \en{should}/\en{must}/\en{have to}),
  12 (thèmes 13 à 16), 17 (partie A : la méthode de lecture).
  \item \textbf{Lecture :} un \en{graded reader} niveau A2 (60 à 80 pages), 10 pages/jour.
  \item \textbf{Vocabulaire :} 1\,500 mots cumulés.
\end{itemize}
\textbf{Jalon :} vous lisez un texte A2 inconnu et en comprenez l'essentiel sans dictionnaire ;
vous parlez de votre métier pendant 5 minutes.

\subsection{Mois 6 --- Bilan de mi-parcours (A2\ra B1)}
\textbf{Priorités :} test complet, conditionnels 0/1/2, propositions relatives,
gérondif contre infinitif, première série de \en{phrasal verbs}, première conversation
libre de 20 minutes.

\begin{itemize}
  \item \textbf{Chapitres :} 10 (unités 4 à 5 : conditionnels 0/1/2, propositions relatives,
  gérondif contre infinitif), 13 (les 75 premiers \emph{phrasal verbs}),
  16 (le protocole des 5 passes), 19 (test du mois 6).
  \item \textbf{Parole :} \textbf{deux conversations libres de 30 minutes par semaine} avec un
  partenaire (application, club, camarade). Sujet imposé : « mon pays, ma famille, mes projets ».
  \item \textbf{Vocabulaire :} 1\,800 mots cumulés.
\end{itemize}
\textbf{Jalon :} test de mi-parcours réussi (chapitre 19) ; vous comprenez une série
avec sous-titres \emph{anglais} ; vous tenez 10 minutes de conversation.

\subsection{Mois 7 --- La nuance (B1)}
\textbf{Priorités :} conditionnel 3, discours rapporté, \en{wish}/\en{if only},
quantifieurs, vocabulaire des médias et de l'actualité, débat.

\begin{itemize}
  \item \textbf{Chapitres :} 11 (unités 1 à 3 : conditionnel 3, discours rapporté,
  \en{wish}/\en{if only}), 13 (expressions idiomatiques 1 à 60), 14 (situations 21 à 28).
  \item \textbf{Écoute :} passage aux podcasts « naturels lents » et aux journaux télévisés.
  \item \textbf{Vocabulaire :} 2\,100 mots cumulés.
\end{itemize}
\textbf{Jalon :} vous donnez et défendez une opinion pendant 10 minutes ; vous comprenez
un journal télévisé à 70\,\%.

\subsection{Mois 8 --- Voyager et présenter (B1)}
\textbf{Priorités :} voix passive complète, \en{have something done}, structures emphatiques,
modaux du passé, situations de voyage (aéroport, hôtel, douane, urgence).

\begin{itemize}
  \item \textbf{Chapitres :} 11 (unités 4 à 7 : passif complet, \en{have something done}, structures emphatiques, modaux du passé), 14 (situations 29 à 36), 16 (compréhension des accents).
  \item \textbf{Projet :} préparer et enregistrer une \textbf{présentation de 5 minutes}
  (votre ville, votre métier, un sujet qui vous passionne).
  \item \textbf{Vocabulaire :} 2\,400 mots cumulés.
\end{itemize}
\textbf{Jalon :} vous gérez seul un trajet complet en anglais (aéroport, hôtel, taxi,
réclamation) ; votre présentation de 5 minutes est fluide.

\subsection{Mois 9 --- Les registres (B1\ra B2)}
\textbf{Priorités :} connecteurs, collocations, idiomes, registre formel/familier,
écriture (e-mails, lettres, CV).

\begin{itemize}
  \item \textbf{Chapitres :} 13 (expressions 61 à 120, argot, les 150 faux amis), 17 (parties B à E : orthographe, ponctuation, les 8 modèles, le journal), annexe D (connecteurs et prépositions).
  \item \textbf{Écriture :} 3 e-mails par semaine (formel, informel, réclamation) + un
  texte de 250 mots chaque dimanche.
  \item \textbf{Vocabulaire :} 2\,700 mots cumulés.
\end{itemize}
\textbf{Jalon :} vous écrivez un e-mail professionnel correct du premier coup ; vous
reconnaissez et utilisez 100 expressions idiomatiques.

\subsection{Mois 10 --- L'anglais professionnel (B2)}
\textbf{Priorités :} entretien d'embauche, réunion, téléphone, négociation, présentation
avec supports.

\begin{itemize}
  \item \textbf{Chapitres :} 14 (situations 37 à 40 et blocs professionnels), 17 (e-mail professionnel, CV, lettre de motivation), 18 (erreurs 31 à 100).
  \item \textbf{Projet :} simuler un entretien d'embauche complet (45 min) avec un
  partenaire ou en auto-questionnement enregistré.
  \item \textbf{Vocabulaire :} 3\,000 mots cumulés.
\end{itemize}
\textbf{Jalon :} vous passez un entretien de 30 minutes ; vous animez une réunion simple.

\subsection{Mois 11 --- La fluidité (B2)}
\textbf{Priorités :} débit, \emph{fillers} naturels, reformulation, humour, compréhension
des accents (britannique, américain, nigérian, ghanéen, indien, sud-africain).

\begin{itemize}
  \item \textbf{Chapitres :} 15 (les 7 drills de fluidité, les \emph{fillers}, le plan
  anti-blocage), 16 (les accents du monde), 6 (révision intonation et \en{connected speech}).
  \item \textbf{Parole :} 30 minutes de conversation libre par jour si possible.
\end{itemize}
\textbf{Jalon :} vous parlez 20 minutes sur un sujet inconnu, sans préparation, sans
repasser au français.

\subsection{Mois 12 --- Certification et autonomie (B2\ra C1)}
\textbf{Priorités :} révision générale, test blanc complet, préparation éventuelle d'un
examen (TOEIC, IELTS, Cambridge B2 First), construction de votre plan pour l'année 2.

\begin{itemize}
  \item \textbf{Chapitres :} 19 (tous les tests et la préparation aux certifications), annexe C (les 14 fiches de révision), annexe F (le bilan des 365 jours).
  \item \textbf{Projet final :} trois conversations d'une heure, un film entier sans
  sous-titres, une présentation de 10 minutes.
\end{itemize}
\textbf{Jalon final :} comparez votre enregistrement du jour 0 à celui du jour 365.
C'est là que vous saurez.

\section{Que faire quand vous prenez du retard}
Vous \emph{allez} prendre du retard : maladie, travail, voyage, mariage, Ramadan, examens,
enfants. C'est normal et ce n'est pas grave --- à une condition : ne jamais transformer un
retard en abandon.

\begin{regle}[La règle des deux jours]
Vous pouvez rater \textbf{un} jour. Jamais \textbf{deux} de suite. Le deuxième jour manqué
est celui qui crée l'habitude de manquer. Si vous ratez deux jours, reprenez le troisième
avec une séance de 15 minutes seulement : l'important est de recoller la chaîne.
\end{regle}

\begin{itemize}
  \item \textbf{Retard de 1 à 2 semaines :} ne « rattrapez » pas en doublant les séances.
  Reprenez le fil, et décalez simplement la fin du plan de deux semaines.
  \item \textbf{Retard d'un mois ou plus :} refaites le test de positionnement, et reprenez
  au mois correspondant, en gardant le vocabulaire SRS à jour (c'est lui qui vous a permis
  de ne rien perdre).
  \item \textbf{Découragement :} relisez votre enregistrement du mois précédent. Les progrès
  sont invisibles au jour le jour, évidents à trois mois.
\end{itemize}

\begin{memo}[Le contrat avec vous-même]
Écrivez et signez, aujourd'hui :\\[4pt]
\textit{Je, \trou[5cm], m'engage à étudier l'anglais \trou[1.5cm] minutes par jour,
\trou[1cm] jours par semaine, du \trou[2.5cm] au \trou[2.5cm], et à parler à voix haute
chaque jour. Mon objectif au \trou[2.5cm] (dans 12 mois) : tenir une conversation
de 30 minutes en anglais sans préparation.}\\[4pt]
Signature : \trou[5cm] \hfill Date : \trou[3cm]
\end{memo}
